\clearpage
\begin{trapbox}
\par\medskip
\noindent\begin{minipage}{\linewidth}
\paragraph{易错一：看到 $f'(c)=0$，就以为原函数一定有两个零点}
\begin{center}% F081: precise mathematical teaching diagram.
\begin{tikzpicture}[x=1cm,y=1cm,>=stealth,font=\small]
\draw[rounded corners=3pt,gray!40] (0,0) rectangle (14.2,8.6);
\begin{scope}[yshift=1cm]
\node at (7.1,7.0) {易错一：都有水平切线，原函数零点却可以是2、1、0个};
\node[] at (2.6,6.3) {$x^2-1$};
\begin{scope}[shift={(2.6,3.6)},x=1.1cm,y=1.0cm]
\draw[->,gray!65] (-1.65,0)--(1.65,0) node[right] {$x$};
\draw[->,gray!65] (0,-1.6)--(0,2) node[left] {$y$};
\draw[blue!75!black,line width=1.15pt,,domain=-1.35:1.35,samples=160] plot (\x,{(\x)*(\x)+(-1)});
\draw[orange!85!black,line width=1.2pt] (-0.32,-1)--(0.32,-1);
\fill[orange!85!black] (0,-1) circle (1.8pt);
\fill[blue!75!black] (-1,0) circle (1.8pt);
\fill[blue!75!black] (1,0) circle (1.8pt);
\end{scope}
\node[] at (2.6,1.15) {2个零点};
\node[] at (7.1,6.3) {$x^2$};
\begin{scope}[shift={(7.1,3.6)},x=1.1cm,y=1.0cm]
\draw[->,gray!65] (-1.65,0)--(1.65,0) node[right] {$x$};
\draw[->,gray!65] (0,-1.6)--(0,2) node[left] {$y$};
\draw[blue!75!black,line width=1.15pt,,domain=-1.35:1.35,samples=160] plot (\x,{(\x)*(\x)+(0)});
\draw[orange!85!black,line width=1.2pt] (-0.32,0)--(0.32,0);
\fill[orange!85!black] (0,0) circle (1.8pt);
\fill[blue!75!black] (0,0) circle (1.8pt);
\end{scope}
\node[] at (7.1,1.15) {1个零点};
\node[] at (11.6,6.3) {$x^2+1$};
\begin{scope}[shift={(11.6,3.6)},x=1.1cm,y=1.0cm]
\draw[->,gray!65] (-1.65,0)--(1.65,0) node[right] {$x$};
\draw[->,gray!65] (0,-1.6)--(0,2) node[left] {$y$};
\draw[blue!75!black,line width=1.15pt,,domain=-1.35:1.35,samples=160] plot (\x,{(\x)*(\x)+(1)});
\draw[orange!85!black,line width=1.2pt] (-0.32,1)--(0.32,1);
\fill[orange!85!black] (0,1) circle (1.8pt);
\end{scope}
\node[] at (11.6,1.15) {0个零点};
\end{scope}
\node[text width=13cm,align=center] at (7.1,.9) {三个函数都满足 $f^{\prime}(0)=0$；存在导数零点不能反推存在两个函数零点。};
\end{tikzpicture}
\end{center}
\end{minipage}\par\medskip

\textbf{错误想法：}导数在某个位置等于 $0$，所以曲线一定会有两次与 $x$ 轴相交。

\textbf{反例：}
\[
f(x)=x^2+1,\qquad f'(0)=0.
\]
但 $x^2+1>0$，原函数一个零点也没有。

\textbf{记住：}\textbf{两个零点能推出中间有导数为零的位置；反过来不能直接推出两个零点。}

\end{trapbox}
\clearpage
\begin{trapbox}
\par\medskip
\noindent\begin{minipage}{\linewidth}
\paragraph{易错二：把“最多三个”写成“恰好三个”}
\begin{center}% F081: precise mathematical teaching diagram.
\begin{tikzpicture}[x=1cm,y=1cm,>=stealth,font=\small]
\draw[rounded corners=3pt,gray!40] (0,0) rectangle (14.2,9.4);
\begin{scope}[yshift=1cm]
\node at (7.1,7.800000000000001) {易错二：导数零点一样多，原函数零点未必一样多};
\node[] at (3.7,7) {$x^3-3x+1$：3个零点};
\begin{scope}[shift={(3.7,3.2)},x=0.92cm,y=0.42cm]
\draw[->,gray!65] (-2.9,0)--(2.6,0) node[right] {$x$};
\draw[->,gray!65] (0,-3.8)--(0,7.2) node[left] {$y$};
\begin{scope}\clip (-2.7,-3.5) rectangle (2.4,7.7);
\draw[blue!75!black,line width=1.15pt,,domain=-2.7:2.4,samples=160] plot (\x,{(\x)*(\x)*(\x)-3*\x+1});
\end{scope}
\fill[blue!75!black] (-1.8793852415718,0) circle (1.8pt);
\fill[blue!75!black] (0.3472963553339,0) circle (1.8pt);
\fill[blue!75!black] (1.532088886238,0) circle (1.8pt);
\draw[orange!85!black,line width=1.2pt] (-1.3,3)--(-0.7,3);
\fill[orange!85!black] (-1,3) circle (1.8pt);
\draw[orange!85!black,line width=1.2pt] (0.7,-1)--(1.3,-1);
\fill[orange!85!black] (1,-1) circle (1.8pt);
\end{scope}
\node[] at (10.4,7) {$x^3-3x+5$：1个零点};
\begin{scope}[shift={(10.4,3.2)},x=0.92cm,y=0.42cm]
\draw[->,gray!65] (-2.9,0)--(2.6,0) node[right] {$x$};
\draw[->,gray!65] (0,-3.8)--(0,7.2) node[left] {$y$};
\begin{scope}\clip (-2.7,-3.5) rectangle (2.4,7.7);
\draw[blue!75!black,line width=1.15pt,,domain=-2.7:2.4,samples=160] plot (\x,{(\x)*(\x)*(\x)-3*\x+5});
\end{scope}
\fill[blue!75!black] (-2.2790187861666,0) circle (1.8pt);
\draw[orange!85!black,line width=1.2pt] (-1.3,7)--(-0.7,7);
\fill[orange!85!black] (-1,7) circle (1.8pt);
\draw[orange!85!black,line width=1.2pt] (0.7,3)--(1.3,3);
\fill[orange!85!black] (1,3) circle (1.8pt);
\end{scope}
\end{scope}
\node[text width=13cm,align=center] at (7.1,.9) {两图同尺度，导数都只有两个零点；只能说最多三个，不能说恰好三个。};
\end{tikzpicture}
\end{center}
\end{minipage}\par\medskip

若 $f'(x)$ 只有两个不同零点，罗尔定理只能保证 $f(x)$ \textbf{最多三个零点}，不能保证确实有三个。

比如
\[
f(x)=x^3-3x+5,
\]
导数同样在 $-1,1$ 处等于 $0$，但 $f(-1)=7>0$、$f(1)=3>0$，整条曲线只有左侧一次穿过 $x$ 轴，实际只有\textbf{一个}零点。

\textbf{正确做法：}上限由导数控制；要说明实际有几个，还需检查连续图像在某些位置是否从负变正或从正变负。

\end{trapbox}
\clearpage
\begin{trapbox}
\par\medskip
\noindent\begin{minipage}{\linewidth}
\paragraph{易错三：只检查两端相等，忘记中间可能有尖角}
\begin{center}% F081: precise mathematical teaching diagram.
\begin{tikzpicture}[x=1cm,y=1cm,>=stealth,font=\small]
\draw[rounded corners=3pt,gray!40] (0,0) rectangle (14.2,8.3);
\begin{scope}[yshift=1cm]
\node at (7.1,6.7) {易错三：连续、两端等高，还必须检查内部可导};
\node[] at (3.8,6) {尖角：内部不可导};
\begin{scope}[shift={(3.8,3.9)},x=1.9cm,y=1.5cm]
\draw[->,gray!65] (-1.35,0)--(1.35,0) node[right] {$x$};
\draw[->,gray!65] (0,-1.45)--(0,0.8) node[above] {$y$};
\draw[blue!75!black,line width=1.15pt,,domain=-1:1,samples=160] plot (\x,{abs(\x)-1});
\fill[blue!75!black] (-1,0) circle (1.8pt);\node[above left] at (-1,0) {$-1$};
\fill[blue!75!black] (1,0) circle (1.8pt);\node[above right] at (1,0) {$1$};
\fill[orange!85!black] (0,-1) circle (1.8pt);
\end{scope}
\node[] at (3.8,1.3) {$|x|-1$};
\node[] at (10.3,6) {圆滑最低点：导数为零};
\begin{scope}[shift={(10.3,3.9)},x=1.9cm,y=1.5cm]
\draw[->,gray!65] (-1.35,0)--(1.35,0) node[right] {$x$};
\draw[->,gray!65] (0,-1.45)--(0,0.8) node[above] {$y$};
\draw[blue!75!black,line width=1.15pt,,domain=-1:1,samples=160] plot (\x,{(\x)*(\x)-1});
\fill[blue!75!black] (-1,0) circle (1.8pt);\node[above left] at (-1,0) {$-1$};
\fill[blue!75!black] (1,0) circle (1.8pt);\node[above right] at (1,0) {$1$};
\fill[orange!85!black] (0,-1) circle (1.8pt);
\draw[orange!85!black,line width=1.2pt] (-0.32,-1)--(0.32,-1);
\end{scope}
\node[] at (10.3,1.3) {$x^2-1$};
\end{scope}
\node[text width=13cm,align=center] at (7.1,.9) {同尺度对照：左图两侧斜率为 $-1,1$，在尖角处接不上同一个导数。};
\end{tikzpicture}
\end{center}
\end{minipage}\par\medskip

看
\[
f(x)=|x|-1,\qquad -1\le x\le1.
\]
两端都有 $f(-1)=f(1)=0$，图像也连着，但在 $x=0$ 是一个尖角，\textbf{那里不能求导}。其他内部位置的导数只可能是 $-1$ 或 $1$，从不等于 $0$。

\textbf{正确做法：}先检查内部每一点是否可导，再使用结论。罗尔定理要求\textbf{中间每一点都能求导}，这里不满足。
\end{trapbox}
\clearpage
\begin{trapbox}
\par\medskip
\noindent\begin{minipage}{\linewidth}
\paragraph{易错四：忽略端点是否连得上}
\begin{center}% F081: precise mathematical teaching diagram.
\begin{tikzpicture}[x=1cm,y=1cm,>=stealth,font=\small]
\draw[rounded corners=3pt,gray!40] (0,0) rectangle (14.2,8.5);
\begin{scope}[yshift=1cm]
\node at (7.1,6.9) {易错四：端点函数值相等，不代表图像接得上};
\node[] at (3.7,6.2) {右端点跳开：不能套用};
\begin{scope}[shift={(2.0,2.8)},x=3.3cm,y=2.2cm]
\draw[->,gray!65] (-0.15,0)--(1.35,0) node[right] {$x$};
\draw[->,gray!65] (0,-0.6)--(0,1.2) node[above] {$y$};
\draw[blue!75!black,line width=1.15pt,,domain=0:1,samples=160] plot (\x,{\x});
\draw[blue!75!black,fill=white,line width=1pt] (1,1) circle (2.2pt);
\fill[blue!75!black] (1,0) circle (1.8pt);\node[below right] at (1,0) {$(1,0)$};
\fill[blue!75!black] (0,0) circle (1.8pt);\node[below left] at (0,0) {$(0,0)$};
\end{scope}
\node[] at (10.2,6.2) {两端连上：可以套用};
\begin{scope}[shift={(8.5,2.8)},x=3.3cm,y=2.2cm]
\draw[->,gray!65] (-0.15,0)--(1.35,0) node[right] {$x$};
\draw[->,gray!65] (0,-0.6)--(0,1.2) node[above] {$y$};
\draw[blue!75!black,line width=1.15pt,,domain=0:1,samples=160] plot (\x,{(\x)*(\x-1)});
\fill[blue!75!black] (1,0) circle (1.8pt);\node[above right] at (1,0) {$(1,0)$};
\draw[orange!85!black,line width=1.2pt] (0.3,-0.25)--(0.7,-0.25);
\fill[blue!75!black] (0,0) circle (1.8pt);\node[below left] at (0,0) {$(0,0)$};
\end{scope}
\end{scope}
\node[text width=13cm,align=center] at (7.1,.9) {两图同尺度；左图 $f(1)=0$，但左侧趋近值为 $1$；右图为 $x(x-1)$。};
\end{tikzpicture}
\end{center}
\end{minipage}\par\medskip

在 $[0,1]$ 上定义
\[
f(x)=\begin{cases}x,&0\le x<1,\\0,&x=1.\end{cases}
\]
虽然 $f(0)=f(1)=0$，且在 $0<x<1$ 时 $f'(x)=1$，但 $x=1$ 处函数值突然从靠近 $1$ 跳到 $0$，\textbf{端点处不连续}。

因此不能套用罗尔定理。\textbf{在整个闭区间连续}和\textbf{中间能求导}，两项都要检查。

\end{trapbox}
\clearpage
\begin{trapbox}
\par\medskip
\noindent\begin{minipage}{\linewidth}
\paragraph{易错五：两个零点不在同一段连续的区间里}
\begin{center}% F081: precise mathematical teaching diagram.
\begin{tikzpicture}[x=1cm,y=1cm,>=stealth,font=\small]
\draw[rounded corners=3pt,gray!40] (0,0) rectangle (14.2,8.5);
\begin{scope}[yshift=1cm]
\node at (7.1,6.9) {易错五：两个零点之间缺了一点};
\begin{scope}[shift={(7.1,3.9)},x=2.2cm,y=1.1cm]
\draw[->,gray!65] (-2.45,0)--(2.45,0) node[right] {$x$};
\draw[->,gray!65] (0,-2)--(0,2) node[above] {$y$};
\draw[blue!75!black,line width=1.15pt,,domain=-2:-0.002,samples=160] plot (\x,{\x+1});
\draw[blue!75!black,line width=1.15pt,,domain=0.002:2,samples=160] plot (\x,{\x-1});
\draw[blue!75!black,fill=white,line width=1pt] (0,1) circle (2.3pt) (0,-1) circle (2.3pt);
\fill[blue!75!black] (-1,0) circle (1.8pt);\node[below] at (-1,0) {$-1$};
\fill[blue!75!black] (1,0) circle (1.8pt);\node[above] at (1,0) {$1$};
\end{scope}
\node[] at (7.1,1.3) {$x=0$ 未定义，不能把两段拼成一个适用区间。};
\end{scope}
\node[text width=13cm,align=center] at (7.1,.9) {两段都满足 $f^{\prime}=1$；连完整区间都没有，不能跨缺口使用罗尔定理。};
\end{tikzpicture}
\end{center}
\end{minipage}\par\medskip

考虑
\[
f(x)=\begin{cases}x+1,&-2\le x<0,\\x-1,&0<x\le2.\end{cases}
\]
$x=-1$ 和 $x=1$ 都是零点，但 $x=0$ 没有定义。虽然在两边各自的区间里导数都是 $1$，也\textbf{不能跨过断开的地方}使用罗尔定理。

\textbf{正确做法：}选定两个零点后，先核对它们之间的每个位置都有定义、连续且内部可导。
\end{trapbox}
\clearpage
\begin{trapbox}
\par\medskip
\noindent\begin{minipage}{\linewidth}
\paragraph{易错六：把相同的解重复计算}
\begin{center}% F081: precise mathematical teaching diagram.
\begin{tikzpicture}[x=1cm,y=1cm,>=stealth,font=\small]
\draw[rounded corners=3pt,gray!40] (0,0) rectangle (14.2,9.0);
\begin{scope}[yshift=1cm]
\node at (7.1,7.4) {易错六：重复因式不会增加零点的位置数};
\begin{scope}[shift={(6.6,3.5)},x=1.8cm,y=0.7cm]
\draw[->,gray!65] (-2.6,0)--(2,0) node[right] {$x$};
\draw[->,gray!65] (0,-2.8)--(0,4.4) node[left] {$y$};
\draw[blue!75!black,line width=1.15pt,,domain=-2.15:1.8,samples=160] plot (\x,{(\x-1)*(\x-1)*(\x+2)});
\fill[blue!75!black] (-2,0) circle (1.8pt);\node[below left] at (-2,0) {$-2$};
\fill[blue!75!black] (1,0) circle (1.8pt);\node[below] at (1,0) {$1$};
\draw[orange!85!black,line width=1.2pt] (0.65,0)--(1.35,0);
\end{scope}
\node[] at (7.1,1.1) {$x=-2$ 穿过横轴；$x=1$ 接触横轴；总共只有两个不同位置。};
\end{scope}
\node[text width=13cm,align=center] at (7.1,.9) {$(x-1)^2(x+2)=0$：数轴上 $x=1$ 只占一个位置。};
\end{tikzpicture}
\end{center}
\end{minipage}\par\medskip

方程
\[
(x-1)^2(x+2)=0
\]
只有 $x=1$ 和 $x=-2$ 两个\textbf{不同的解}。虽然因式里 $(x-1)$ 出现两次，但在数轴上 $x=1$ 还是只有一个位置。

\textbf{零点个数关系数的是“不同位置”，不是某个因式出现多少次。}

\end{trapbox}
\clearpage
\begin{trapbox}
\par\medskip
\noindent\begin{minipage}{\linewidth}
\paragraph{易错七：把 $c$ 当成两端的正中点，或把端点算进去}
\begin{center}% F081: precise mathematical teaching diagram.
\begin{tikzpicture}[x=1cm,y=1cm,>=stealth,font=\small]
\draw[rounded corners=3pt,gray!40] (0,0) rectangle (14.2,9.0);
\begin{scope}[yshift=1cm]
\node at (7.1,7.4) {易错七：内部的水平切线不必位于正中间};
\begin{scope}[shift={(3.2,5.4)},x=7cm,y=16cm]
\draw[->,gray!65] (-0.08,0)--(1.2,0) node[right] {$x$};
\draw[->,gray!65] (0,-0.19)--(0,0.06) node[above] {$y$};
\draw[blue!75!black,line width=1.15pt,,domain=0:1,samples=160] plot (\x,{(\x)*(\x)*(\x-1)});
\fill[blue!75!black] (0,0) circle (1.8pt);\node[above left] at (0,0) {$a=0$};
\fill[blue!75!black] (1,0) circle (1.8pt);\node[above right] at (1,0) {$b=1$};
\fill[orange!85!black] (0.6666666666666666,-0.14814814814814814) circle (1.8pt);
\draw[orange!85!black,line width=1.2pt] (0.5266666666666666,-0.14814814814814814)--(0.8066666666666666,-0.14814814814814814);
\draw[dashed,gray] (.5,0)--(.5,-.16);\draw[dashed,orange!85!black] (.6666667,0)--(.6666667,-.1481481);
\node[above] at (.5,0) {$\frac12$};\node[below] at (.6666667,-.155) {$c=\frac23$};
\end{scope}
\node[] at (7.1,1.25) {$f^{\prime}(x)=x(3x-2)$；端点 $0$ 的导数也为零，但不能作为内部的 $c$。};
\end{scope}
\node[text width=13cm,align=center] at (7.1,.9) {$f(x)=x^2(x-1)$ 在 $[0,1]$ 满足罗尔条件，内部位置为 $c=2/3$。};
\end{tikzpicture}
\end{center}
\end{minipage}\par\medskip

罗尔定理只保证存在
\[
\boxed{\ a<c<b,\qquad f'(c)=0\ }.
\]
\textbf{它没说 $c=(a+b)/2$，也没说可以取 $c=a$ 或 $c=b$。}对称的二次函数会在中点取得这个位置；一般函数须依据导数求解，不能预先指定。图中 $f(x)=x^2(x-1)$ 的内部位置为 $2/3$，不是中点 $1/2$；虽然端点 $0$ 的导数也为零，它仍不属于所要求的开区间。
\end{trapbox}

\clearpage
\begin{trapbox}
\paragraph{易错八：认为水平切线一定出现在最高处或最低处}

\textbf{错误想法：}曲线只有先升后降、或先降后升时，才能有水平切线。

看 $f(x)=x^3$，它从左到右始终增大，但
\[
 f'(x)=3x^2,\qquad f'(0)=0.
\]
原点处切线就是水平的，图像却没有在这里从上升变为下降。

\textbf{正确理解：}罗尔定理的\textbf{结论只要求切线水平}；它并不承诺每条水平切线都位于最高点或最低点。
\end{trapbox}
\clearpage
\begin{trapbox}
\paragraph{易错九：认为水平切线一定与两端的连线重合}

取 $f(x)=(x-2)^2+1$，有
\[
 A(1,2),\quad B(3,2),\quad C(2,1).
\]
直线 $AB$ 是 $y=2$，而点 $C$ 处的水平切线是 $y=1$。\textbf{它们平行，却不是同一条直线}。

\textbf{正确理解：}两端相同的是\textbf{高度}；中间找到的是\textbf{切线方向}。不能把“切线与两端连线平行”误读成“切点一定与两端同高”。
\end{trapbox}


% -------- 新增易错：九种是有条件的基本分类 --------
\clearpage
\begin{trapbox}
\paragraph{易错十：把九种走向当成所有函数的全部情况}

\textbf{误区一：}看到 $f'(c)=0$，就把它当成``左右两边导数都等于零''。
实际上这只说明\textbf{在 $c$ 这一点}导数是 $0$。表中某一侧写``平''，要求\textbf{那一侧的一整小段导数都为 $0$}。

\textbf{误区二：}以为九宫格可以不加条件地套在所有函数上。
九种组合的前提是：\textbf{左右每一侧紧邻 $c$ 的一小段里，导数的符号各自固定不变}。
若一侧还有反复上升、下降的变化，就先把那侧继续分段，不能硬归成一个符号。

\textbf{误区三：}把``一处水平切线''当作``一定有最高点或最低点''。
例如 $f(x)=x^3$ 在 $0$ 处的切线水平，但左右两侧的导数均为正，\textbf{图像并未掉头}。

\textbf{答题顺序：}\textbf{先查是否可导 $\to$ 看左右导数符号 $\to$ 判断是否掉头}；
不要只凭 $f'(c)=0$ 就判断曲线走向。
\end{trapbox}
